\documentclass[11pt,letterpaper,oneside]{report}
\usepackage[T1]{fontenc}
\usepackage[utf8]{inputenc}
\usepackage{lmodern}
\usepackage[margin=0.90in,headheight=15pt]{geometry}
\usepackage{microtype}
\usepackage[svgnames,table]{xcolor}
\usepackage{amsmath,amssymb,booktabs,longtable,tabularx,array}
\usepackage{enumitem}
\usepackage[most]{tcolorbox}
\usepackage{listings}
\usepackage{fancyhdr}
\usepackage{titlesec}
\usepackage{hyperref}
\usepackage{xurl}
\usepackage{lastpage}
\definecolor{ink}{HTML}{14273C}
\definecolor{primary}{HTML}{0B5379}
\definecolor{accent}{HTML}{B66A30}
\definecolor{pale}{HTML}{F1F5F8}
\definecolor{codebg}{HTML}{F5F7F8}
\hypersetup{colorlinks=true,linkcolor=primary,urlcolor=primary,pdftitle={Linux and UNIX System Programming Engineering Suite}}
\titleformat{\chapter}[display]{\normalfont\LARGE\bfseries\color{ink}}{\color{accent}\large\MakeUppercase{\chaptername}\ \thechapter}{10pt}{}[\vspace{0.4ex}\titlerule]
\titleformat{\section}{\Large\bfseries\color{primary}}{\thesection}{0.7em}{}
\titleformat{\subsection}{\large\bfseries\color{ink}}{\thesubsection}{0.7em}{}
\setlist[itemize]{leftmargin=*,topsep=3pt,itemsep=2pt}
\setlist[enumerate]{leftmargin=*,topsep=3pt,itemsep=3pt}
\setlength{\parskip}{5pt}
\setlength{\parindent}{0pt}
\lstdefinestyle{sysc}{language=C,basicstyle=\ttfamily\footnotesize,backgroundcolor=\color{codebg},keywordstyle=\bfseries\color{primary},commentstyle=\color{ForestGreen},stringstyle=\color{Brown},breaklines=true,showstringspaces=false,tabsize=4,frame=single,rulecolor=\color{LightSteelBlue},columns=fullflexible,keepspaces=true}
\lstset{style=sysc}
\newtcolorbox{engineeringnote}[1][]{colback=pale,colframe=primary,boxrule=0.55pt,arc=1.5mm,title=Engineering note,fonttitle=\bfseries,#1}
\newtcolorbox{modernnote}[1][]{colback=white,colframe=accent,boxrule=0.55pt,arc=1.5mm,title=Modern extension --- not in the 2010 textbook,fonttitle=\bfseries,#1}
\newcommand{\sys}[1]{\texttt{#1}}
\pagestyle{fancy}\fancyhf{}\fancyhead[L]{\small\color{ink} Linux / UNIX Engineering Suite}\fancyhead[R]{\small\color{primary}\nouppercase{\leftmark}}\fancyfoot[C]{\small\thepage\ / \pageref{LastPage}}
\renewcommand{\headrulewidth}{0.45pt}
\setcounter{tocdepth}{1}
\setlength{\emergencystretch}{2em}
\begin{document}
\begin{titlepage}\centering\vspace*{2cm}{\color{primary}\Large ENGINEERING REFERENCE SERIES\par}\vspace{1.25cm}{\color{ink}\Huge\bfseries Concurrency and IPC Handbook\par}\vspace{0.6cm}{\Large Pthreads, processes, synchronization and message transport\par}\vspace{2cm}\rule{0.75\textwidth}{1.2pt}\par\vspace{0.55cm}{\large An original companion to Michael Kerrisk\textquotesingle s 2010 text}\par\vspace{0.3cm}{\large Edition 1.0\quad |\quad October 2026}\par\vfill{\small Document 06\quad |\quad pdfLaTeX-compatible standalone source}\par\end{titlepage}
\tableofcontents\clearpage
\chapter{Ownership and progress model}
Concurrency control requires two distinct arguments: safety, meaning forbidden states never occur; and liveness, meaning work can eventually advance under stated scheduling and availability assumptions. A lock prevents certain races but can create deadlocks or starvation. A queue provides a communication abstraction but introduces capacity and shutdown state.

\section{Progress properties}
\begin{itemize}
\item Mutual exclusion: only one thread can mutate a protected invariant at a time.
\item Deadlock freedom: no cycle of permanent waits on locks, queues or other resources under declared assumptions.
\item Starvation resistance: a participant is not indefinitely excluded under intended load.
\item Backpressure: memory and queued work stay within explicit bounds.
\item Termination: waiting threads and subprocesses can observe a closed state and exit.
\end{itemize}
\chapter{Mutex and condition variable protocol}
The guarded predicate is the source of truth. A signal or broadcast only announces that rechecking may succeed. Condition-variable waits release and reacquire a mutex atomically, but wakeups may occur without the application predicate becoming true.

\begin{lstlisting}
#define _POSIX_C_SOURCE 200809L
#include <pthread.h>
#include <stdio.h>
#include <stdlib.h>

#define N 8
struct queue {
    int a[N], head, tail, count, closed;
    pthread_mutex_t mu;
    pthread_cond_t readable, writable;
};
static struct queue q = { .mu = PTHREAD_MUTEX_INITIALIZER,
    .readable = PTHREAD_COND_INITIALIZER,
    .writable = PTHREAD_COND_INITIALIZER };
static void *consumer(void *unused) {
    (void)unused;
    long sum = 0;
    pthread_mutex_lock(&q.mu);
    for (;;) {
        while (q.count == 0 && !q.closed)
            pthread_cond_wait(&q.readable, &q.mu);
        if (q.count == 0 && q.closed) break;
        int x = q.a[q.head]; q.head = (q.head + 1) % N; --q.count;
        pthread_cond_signal(&q.writable);
        pthread_mutex_unlock(&q.mu);
        sum += x;
        pthread_mutex_lock(&q.mu);
    }
    pthread_mutex_unlock(&q.mu);
    printf("sum=%ld\n", sum);
    return NULL;
}
int main(void) {
    pthread_t t;
    int rc = pthread_create(&t, NULL, consumer, NULL);
    if (rc != 0) { fprintf(stderr, "pthread_create: %d\n", rc); return 1; }
    for (int i = 1; i <= 100; ++i) {
        pthread_mutex_lock(&q.mu);
        while (q.count == N)
            pthread_cond_wait(&q.writable, &q.mu);
        q.a[q.tail] = i; q.tail = (q.tail + 1) % N; ++q.count;
        pthread_cond_signal(&q.readable);
        pthread_mutex_unlock(&q.mu);
    }
    pthread_mutex_lock(&q.mu);
    q.closed = 1;
    pthread_cond_broadcast(&q.readable);
    pthread_mutex_unlock(&q.mu);
    rc = pthread_join(t, NULL);
    if (rc != 0) { fprintf(stderr, "pthread_join: %d\n", rc); return 1; }
    pthread_cond_destroy(&q.writable); pthread_cond_destroy(&q.readable);
    pthread_mutex_destroy(&q.mu);
    return 0;
}
\end{lstlisting}
\section{State transition proof}
\begin{enumerate}
\item Invariant I1: 0 <= count <= N.
\item Invariant I2: head and tail are updated only while mutex mu is held.
\item Producer waits while count == N and observes the room predicate under the same mutex.
\item Consumer waits while count == 0 and closed is false.
\item Setting closed and broadcasting allows a blocked consumer to terminate once the queue drains.
\item All heap/stack objects must remain valid through pthread\_join, and the mutex/conditions are destroyed last.
\end{enumerate}
\chapter{Deadlock prevention and cancellation}
\begin{itemize}
\item Maintain a global acquisition order across locks.
\item Do not block on external I/O while holding a coarse shared-state mutex.
\item Avoid callbacks into arbitrary user code while holding internal locks.
\item Use deferred cancellation and cleanup handlers around protected resources.
\item Use deadline-aware waits, and treat timed waits as potentially racing with normal completion.
\end{itemize}
\chapter{Process communication choices}
\begin{longtable}{p{.19\linewidth}p{.26\linewidth}p{.22\linewidth}p{.22\linewidth}}\toprule
Facility & Transfer unit & Coordination & Failure/lifetime concern\\\midrule\endhead
pipe / FIFO & bytes & Blocking/nonblocking read/write & EOF and surviving writers\\
UNIX socket & stream or datagram & send/recv with peer model & Connection shutdown and credentials\\
POSIX mq & whole message + priority & queue management & Message limit and name cleanup\\
SysV mq & typed messages & ID and key-based access & Persistent kernel object cleanup\\
POSIX shm & shared bytes & Separate mutex/semaphore & ABI layout and initialization\\
SysV shm & shared bytes & Separate sync primitive & Segment lifetime and attach semantics\\
POSIX semaphore & integer count & \sys{sem\_wait}/\sys{sem\_post} & No owner recovery guarantee\\
fcntl locks & byte ranges & Cooperative file participants & Lock ownership and advisory model\\
\bottomrule\end{longtable}\chapter{IPC protocol specification template}
\section{Binary header contract}
For any IPC mechanism that transports untrusted messages, define a version, message type, fixed header size, payload size, byte order, integrity policy, correlation token and timeout rule. Never cast an untrusted byte buffer directly to a padding-sensitive C struct unless the ABI is explicitly controlled.

\begin{center}\fcolorbox{primary}{pale}{\parbox{.90\linewidth}{\texttt{magic:4 | version:2 | type:2 | length:4 | request-id:8 | payload:length}\par\smallskip
All fields have declared network byte order; reject length above a configured maximum before allocation.}}
\end{center}\section{State machine}
\begin{enumerate}
\item DISCONNECTED: obtain transport and authenticate endpoint.
\item NEGOTIATING: check version and capabilities, with size and time bounds.
\item READY: accept frames, validate sizes and enqueue bounded work.
\item DRAINING: reject new work, finish or cancel in-flight requests, release resources.
\item CLOSED: close handles, unlink owned objects and record termination result.
\end{enumerate}
\chapter{Shared memory publication}
A mapped shared region can have different virtual addresses in each process. Use offsets from a versioned region base, never raw process pointers. Synchronization primitives must be initialized for process sharing where the standard supports them. Initialization, recovery after crash, endian/ABI assumptions and writer/reader visibility all require a protocol.

\begin{itemize}
\item Header: magic, version, total length, state, data offsets and optional generation counter.
\item Publisher: validate size, allocate and map, initialize lock and metadata, then mark READY through synchronized publication.
\item Reader: map, verify header and capabilities, acquire synchronization and bounds-check every offset.
\item Teardown: mark DRAINING/CLOSED, stop new attachers, release mappings, unlink owned names.
\item Recovery: define whether an owner dying mid-update can be detected and repaired.
\end{itemize}
\chapter{Signals versus synchronous coordination}
Signals are useful for lifecycle notifications, but standard signals can coalesce and handler execution is constrained. Prefer queued messages for durable work requests, mutexes and condition variables for shared-memory state, and dedicated sigwait threads or signalfd-based event loops where appropriate. Linux signalfd is an explicit extension.

\chapter{Linux eventfd, timerfd, signalfd and epoll}
Linux exposes several synchronization and notification events through file descriptors, enabling a unified readiness loop. This can simplify architecture when an application already uses epoll, but each descriptor has separate consumption semantics and portability limits.

\begin{lstlisting}
#define _GNU_SOURCE
#include <errno.h>
#include <stdint.h>
#include <stdio.h>
#include <stdlib.h>
#include <sys/epoll.h>
#include <sys/eventfd.h>
#include <unistd.h>

int main(void) {
    int e = eventfd(0, EFD_CLOEXEC | EFD_NONBLOCK);
    if (e == -1) { perror("eventfd"); return 1; }
    int ep = epoll_create1(EPOLL_CLOEXEC);
    if (ep == -1) { perror("epoll_create1"); close(e); return 1; }
    struct epoll_event event = { .events = EPOLLIN,
                                 .data.fd = e };
    if (epoll_ctl(ep, EPOLL_CTL_ADD, e, &event) == -1) {
        perror("epoll_ctl"); close(ep); close(e); return 1;
    }
    uint64_t one = 1, got = 0;
    if (write(e, &one, sizeof one) != (ssize_t)sizeof one) {
        perror("write eventfd"); close(ep); close(e); return 1;
    }
    struct epoll_event ready;
    int count = epoll_wait(ep, &ready, 1, 1000);
    if (count != 1 || !(ready.events & EPOLLIN) ||
        read(e, &got, sizeof got) != (ssize_t)sizeof got) {
        fprintf(stderr, "readiness test failed\n");
        close(ep); close(e); return 1;
    }
    printf("event count=%llu\n", (unsigned long long)got);
    close(ep); close(e); return 0;
}
\end{lstlisting}
\chapter{Performance and correctness experiment design}
\begin{longtable}{p{.28\linewidth}p{.59\linewidth}}\toprule
Measured dimension & Measurement discipline\\\midrule\endhead
Throughput & Fixed payload distribution; warm-up, repeated runs, stated CPU count\\
Latency & p50/p95/p99, maximum, arrival process and end-to-end definition\\
CPU & user/system time, context switches and idle baseline\\
Memory & peak RSS, queue bytes and shared-region size\\
Fairness & per-worker completed jobs and starvation indicators\\
Reliability & fault-injected missed wakeups, crash recovery and cleanup\\
\bottomrule\end{longtable}\chapter{Concurrency review gates}
\begin{enumerate}
\item Can every shared variable be associated with a lock, atomic or immutable publication rule?
\item Is there a unique shutdown owner and a finite wait path for every worker?
\item Can lost notifications, spurious wakeups or a peer crash violate safety?
\item Does a queue have explicit size and overload behavior?
\item Does IPC object cleanup work after both graceful and forced termination?
\item Were ThreadSanitizer and stress loops run where compatible with the chosen IPC model?
\end{enumerate}

\chapter*{Source and currency notes}
\addcontentsline{toc}{chapter}{Source and currency notes}
\textbf{Foundational source:} Michael Kerrisk, \emph{The Linux Programming Interface: A Linux and UNIX System Programming Handbook}, No Starch Press, 2010, ISBN 978-1-59327-220-3. The textbook's chapter titles and start pages organize the study guide; explanations, tasks, and code in this suite are independently authored educational material, not reproductions of the original text or its example programs.

\textbf{Current reference baseline:} The Open Group / IEEE, \emph{POSIX.1-2024, Issue 8}, \url{https://pubs.opengroup.org/onlinepubs/9799919799/}; Linux man-pages project, \url{https://man7.org/linux/man-pages/}; Linux kernel documentation, \url{https://docs.kernel.org/}. For Linux-only extensions consult the relevant per-interface manual pages and kernel documentation for the target kernel, architecture and libc; availability can differ materially.

\textbf{Scope:} This is an engineering study and implementation companion. It does not claim to restate a normative standard or guarantee a specific Linux kernel's behavior. System calls, libc wrappers and higher-level library functions are explicitly distinguished where material.
\end{document}
